\section{Results}
\label{sec:results}

Table~\ref{tab:p2scores} is the decision path. Lower log loss is better. The
local numbers are the equal average of the three phase losses on the shared
checks. File~13 is the submitted model. It is the best reported Kaggle score.
File~15 won the local time-ordered average and lost on Kaggle.

\begin{table}[t]
\centering\small
\setlength{\tabcolsep}{3pt}
\begin{tabular*}{\columnwidth}{@{\extracolsep{\fill}}lccc@{}}
\toprule
Model & Time & Grouped & Kaggle \\
\midrule
Phase-specific ensemble & --- & --- & 0.59562 \\
Text-heavy ensemble & --- & --- & 0.58993 \\
Stack, no recency weights & 0.5421 & 0.5797 & 0.58810 \\
Recency stack & 0.5374 & 0.5874 & 0.58688 \\
File 10, recent neighbors & 0.5359 & 0.5869 & 0.58609 \\
File 11, 20\% early ensemble & 0.5339 & 0.5812 & 0.58474 \\
\textbf{File 13, 30\% early ensemble} & 0.5334 & 0.5790 & \textbf{0.58465} \\
File 15, tree and trend added & \textbf{0.5327} & \textbf{0.5783} & 0.58492 \\
\bottomrule
\end{tabular*}
\caption{Later models on the shared checks. Time is the average of the
original forward windows. Grouped is repeated trial-grouped cross-validation.
The first two Kaggle scores are early submissions that used a different local
protocol, so those two local cells are left blank. Kaggle scores are the
reported public scores.}
\label{tab:p2scores}
\end{table}

Table~\ref{tab:p2phase} is the supplied validation split, under preprocessing
fit only on the training rows of that split. The baseline is the criteria
logistic model. The recency stack is the 0.58688 reference. File~10 is the
neighbor blend before the older ensemble is added. File~10 does not win this
table. Its gain showed up on the later-year windows and on Kaggle, which is
why we did not select on validation F1.

\begin{table}[t]
\centering\small
\setlength{\tabcolsep}{2.5pt}
\begin{tabular}{llccc}
\toprule
Phase & Metric & Baseline & Recency & File 10 \\
\midrule
\multirow{4}{*}{I}
 & F1       & 0.765 & \textbf{0.793} & 0.762 \\
 & ROC-AUC  & 0.783 & \textbf{0.794} & 0.787 \\
 & PR-AUC   & \textbf{0.835} & 0.830 & 0.824 \\
 & log loss & 0.556 & 0.556 & 0.560 \\
\midrule
\multirow{4}{*}{II}
 & F1       & \textbf{0.657} & 0.647 & 0.638 \\
 & ROC-AUC  & 0.740 & 0.753 & \textbf{0.754} \\
 & PR-AUC   & 0.717 & \textbf{0.732} & 0.730 \\
 & log loss & 0.600 & 0.593 & 0.593 \\
\midrule
\multirow{4}{*}{III}
 & F1       & 0.833 & 0.825 & \textbf{0.829} \\
 & ROC-AUC  & 0.774 & \textbf{0.790} & 0.786 \\
 & PR-AUC   & 0.862 & \textbf{0.882} & 0.877 \\
 & log loss & 0.529 & \textbf{0.516} & 0.521 \\
\bottomrule
\end{tabular}
\caption{Validation metrics by phase. F1 uses a cutoff of 0.5. ROC-AUC,
PR-AUC, and log loss use the predicted probability that the trial succeeds.
Phase sizes are 149, 501, and 363 rows. Best value per row in bold. Phase~I
log loss ties at three decimals; the recency stack is slightly lower.}
\label{tab:p2phase}
\end{table}

Table~\ref{tab:p2err} is the reason check on the failed validation trials.
The baseline is the criteria logistic model. The ensemble is the early
text-and-tree model that later became the 30\% share of file~13. Most reason
tags are Part~1 predictions, not confirmed labels. The groups large enough
to read are Phase~II Recruitment and the three no-reason groups. Phase~III
failures are called successful far more often than Phase~II failures. Adding
clinical context lowers some of those rates and raises others. The later
neighbor model did not fix the hardest cell: all 9 Phase~III efficacy-tagged
failures were still called successful.

\begin{table}[t]
\centering\small
\setlength{\tabcolsep}{3pt}
\begin{tabular*}{\columnwidth}{@{\extracolsep{\fill}}llccc@{}}
\toprule
Phase & Reason & $n$ & Base & Ens. \\
\midrule
I & Safety & 9 & .56 & .44 \\
I & Funding & 4 & .50 & .50 \\
I & Efficacy & 2 & .50 & .50 \\
I & Recruitment & 22 & .36 & .27 \\
I & Other & 5 & .80 & .80 \\
I & No reason & 28 & .43 & .39 \\
II & Safety & 15 & .40 & .40 \\
II & Funding & 4 & .00 & .25 \\
II & Efficacy & 10 & .30 & .30 \\
II & Recruitment & 87 & .20 & .16 \\
II & Other & 18 & .33 & .28 \\
II & No reason & 144 & .35 & .34 \\
III & Safety & 6 & .50 & .67 \\
III & Funding & 0 & --- & --- \\
III & Efficacy & 9 & .89 & .67 \\
III & Recruitment & 21 & .52 & .57 \\
III & Other & 8 & .63 & .75 \\
III & No reason & 82 & .56 & .55 \\
\bottomrule
\end{tabular*}
\caption{Share of failed validation trials predicted successful, by phase and
stopping reason. ``No reason'' means the trial never appears in Part~1.
A trial with several reasons is listed in each. Funding has no Phase~III
failures.}
\label{tab:p2err}
\end{table}
